
Iterative code repair --- applying feedback signals to guide LLM re-generation on failing problems --- improves functional correctness, but the choice of feedback signal remains underspecified in practice. This paper compares execution test feedback against pylint/mypy static analysis feedback for Llama 3.1 8B Instruct in an iso-compute setting (B=1000 output tokens per problem) on HumanEval (164 problems) and MBPP (378 problems). Execution feedback significantly outperforms pylint/mypy on both benchmarks (McNemar's test: HumanEval p=0.0001, MBPP p<$10^{-18})$, yielding $\Delta _pass$@1 of +4.9pp on HumanEval and +40.2pp on MBPP, while pylint/mypy repair reduces HumanEval pass@1 by 4.3pp and improves MBPP by 18.3pp. The mechanism underlying this divergence is a style-function dissociation: pylint flags 100\% of HumanEval baseline failures, but 94.3\% of those flags correspond to Convention-category style rules (C0304: missing final newline, C0114: missing module docstring) that fire universally on LLM-generated code regardless of functional correctness; functional Error/Warning coverage is 12.5\% across 64 failures. Mypy detects 0\% of failures. These findings suggest that feedback signal selection for LLM code repair should prioritize functional informativeness over aggregate coverage metrics.

